% fig3-routing-results.tex — §7 routing evaluation bar chart (RC-07)
% Data source: data/aggregated/routing-baseline-results.csv
% Columns: router,accuracy,macro_f1,haiku_f1,sonnet_f1,opus_f1,total_cost_saved_usd,total_cost_of_mistakes_usd
% Style: print-safe, grouped bars for accuracy/macro_f1/haiku_f1/sonnet_f1, width=0.85\textwidth

\begin{figure}[ht]
\centering
\pgfplotstableread[col sep=comma]{data/aggregated/routing-baseline-results.csv}{\routingtable}

\begin{tikzpicture}
\begin{axis}[
  ybar,
  width=0.85\textwidth,
  height=6cm,
  bar width=5pt,
  ylabel={Score},
  ymin=0, ymax=1.05,
  ytick={0,0.2,0.4,0.6,0.8,1.0},
  yticklabel style={font=\small},
  xtick=data,
  xticklabels={Always-opus, Random, Length, \textbf{Semantic-sim}},
  xticklabel style={font=\small, align=center},
  enlarge x limits=0.15,
  legend style={
    at={(0.97,0.97)},
    anchor=north east,
    legend columns=2,
    font=\small,
    draw=black,
  },
  nodes near coords,
  nodes near coords style={font=\tiny, rotate=90, anchor=west},
  every node near coord/.append style={/pgf/number format/.cd, fixed, precision=2},
]
% Accuracy
\addplot+[fill=black!80, draw=black] table[x expr=\coordindex, y=accuracy] {\routingtable};
% Macro F1
\addplot+[fill=black!50, draw=black] table[x expr=\coordindex, y=macro_f1] {\routingtable};
% Haiku F1
\addplot+[fill=black!20, draw=black] table[x expr=\coordindex, y=haiku_f1] {\routingtable};
% Sonnet F1
\addplot+[fill=white, draw=black, pattern=north east lines] table[x expr=\coordindex, y=sonnet_f1] {\routingtable};

\legend{Accuracy, Macro-F1, Haiku-F1, Sonnet-F1}
\end{axis}
\end{tikzpicture}

\caption{%
  Routing evaluation on N\,=\,49 held-out turns (LOO-CV for semantic-similarity).
  The proposed semantic-similarity router achieves 61.2\,\% accuracy
  (macro-F1\,=\,0.526), with strong haiku-F1 (0.786) and solid sonnet-F1 (0.667),
  saving \$0.000289 per 49-turn batch vs.\ \$0 for always-opus by construction.
  Source: \texttt{data/aggregated/routing-baseline-results.csv}.%
}
\label{fig:routing-results}
\end{figure}
